\documentclass{../templateWriting/cdwriting}

\newcommand{\authorinfo}{THE FORUM CENTER - NGUYỄN HOÀNG HUY}
\newcommand{\documenttitle}{TÀI LIỆU CHUYÊN ĐỀ WRITING}
\newcommand{\collectiontitle}{Level 1 — W4 (Jan 19-25)}

\begin{document}

\thispagestyle{firstpage}
\makeworkshopheader{\authorinfo}{\documenttitle}{\collectiontitle}
\pagestyle{otherpages}

\workshopprompt{Some people do not go directly to college but travel or work for a short time. Do you think the advantages outweigh the disadvantages?}

\cdsection{Phần I: Xây dựng nền tảng tư duy và chiến lược}

Trước khi viết, chúng ta cần hiểu rõ đề bài và sắp xếp ý tưởng để không bị rối. Với Band 4.0-4.5, sự rõ ràng là quan trọng nhất.

\cdstep{1. Bước 1: Giải mã Đề bài \textit{(Decoding the Prompt)}}
\begin{itemize}
\item \cdkv{Chủ đề cốt lõi}{Việc học sinh không vào đại học ngay mà dành thời gian đi du lịch hoặc làm việc (thường gọi là "Gap Year").}
\item \cdkv{Yêu cầu nhiệm vụ}{Câu hỏi "Do you think the advantages outweigh the disadvantages?" yêu cầu bạn so sánh Lợi ích và Tác hại. Bạn cần quyết định xem cái nào lớn hơn.}
\item \cdkv{Định hướng chiến lược}{Với trình độ này, bạn nên chọn quan điểm rõ ràng: "Lợi ích nhiều hơn Tác hại". Điều này giúp bạn dễ viết hơn vì có nhiều ý tưởng tích cực để nói về trải nghiệm và kiếm tiền.}
\end{itemize}

\cdstep{2. Bước 2: Quy trình Tìm ý \textit{(Brainstorming)}}
\begin{itemize}
\item \cdkv{Xác định Từ khóa}{College/University (Đại học), Work (Làm việc), Travel (Du lịch), Experience (Kinh nghiệm), Money (Tiền).}
\item \cdblue{\textbf{Phát triển Ý tưởng:}}
\begin{itemize}
\item \cdkv{Lợi ích (Advantages)}{Kiếm tiền đóng học phí, học kỹ năng mềm (giao tiếp, tự lập), xả hơi sau 12 năm học phổ thông.}
\item \cdkv{Tác hại (Disadvantages)}{Quên kiến thức, ra trường muộn hơn bạn bè, tốn tiền (nếu đi du lịch nhiều).}
\end{itemize}
\item \cdkv{Sàng lọc Luận điểm}{Bài viết sẽ tập trung vào 2 ý chính:}
\begin{itemize}
\item \cdkv{Tác hại}{Có thể làm mất thói quen học tập.}
\item \cdkv{Lợi ích (Quan trọng hơn)}{Giúp trưởng thành hơn và hiểu giá trị đồng tiền.}
\end{itemize}
\end{itemize}

Kĩ thuật mở rộng ý tưởng (Idea Development tips)

Để viết câu dài hơn và hay hơn, hãy dùng các cách sau:
\begin{itemize}
\item \cdblue{\textbf{Phương pháp So sánh (Contrast):}}
\begin{itemize}
\item \cdkv{Ý tưởng}{Đi làm giúp trưởng thành hơn.}
\item \cdkv{Viết}{Học sinh ở trường chỉ biết lý thuyết, NHƯNG học sinh đi làm sẽ hiểu thực tế cuộc sống.}
\end{itemize}
\item \cdblue{\textbf{Liên hệ Thực tế (Real-world Connection):}}
\begin{itemize}
\item \cdkv{Ý tưởng}{Kiếm tiền là tốt.}
\item \cdkv{Viết}{Nhiều sinh viên nghèo phải đi làm 1 năm để tiết kiệm tiền trả học phí đại học.}
\end{itemize}
\item \cdblue{\textbf{Giả định "Nếu - Thì" (Hypothetical Scenarios):}}
\begin{itemize}
\item \cdkv{Ý tưởng}{Nếu không đi làm trước.}
\item \cdkv{Viết}{Nếu sinh viên vào đại học ngay lập tức, họ có thể cảm thấy mệt mỏi và không biết mình thực sự thích gì.}
\end{itemize}
\end{itemize}

\cdgreenheading{A. Phân tích Tác hại (Disadvantages - Yếu tố tiêu cực)}

Đây là phần chúng ta thừa nhận nhưng sẽ cho là nhỏ hơn lợi ích.

\cdstep{1. Mất thói quen học tập \textit{(Loss of study habits)}}
\begin{itemize}
\item \cdkv{Lập luận}{Sau một năm chỉ làm việc hoặc chơi, quay lại bàn học sẽ rất khó khăn.}
\item \cdkv{Dẫn chứng/Giải thích}{Học sinh có thể thấy sách vở nhàm chán so với việc kiếm ra tiền hoặc đi chơi. Điều này làm giảm động lực học tập.}
\end{itemize}

\cdstep{2. Ra trường muộn hơn \textit{(Delay in graduation)}}
\begin{itemize}
\item \cdkv{Lập luận}{Bạn sẽ tốt nghiệp và đi làm muộn hơn bạn bè cùng trang lứa 1 năm.}
\item \cdkv{Dẫn chứng/Giải thích}{Khi bạn bè đã có công việc ổn định, bạn mới bắt đầu ra trường xin việc. Điều này có thể tạo áp lực tâm lý.}
\end{itemize}

\cdgreenheading{B. Phân tích Lợi ích (Advantages - Yếu tố tích cực)}

Đây là phần trọng tâm để bảo vệ quan điểm của bạn.

\cdstep{1. Trưởng thành và Kỹ năng sống \textit{(Maturity and Life skills)}}
\begin{itemize}
\item \cdkv{Lập luận}{Đi làm hoặc đi xa giúp học sinh tự lập, không dựa dẫm vào bố mẹ.}
\item \cdblue{\textbf{Dẫn chứng/Giải thích:}}
\begin{itemize}
\item \cdkv{Ví dụ}{Khi đi làm thêm (bồi bàn, bán hàng), bạn học được cách giao tiếp với khách hàng và quản lý thời gian. Đây là "soft skills" (kỹ năng mềm) mà trường học ít dạy.}
\end{itemize}
\end{itemize}

\cdstep{2. Hiểu rõ định hướng tương lai \textit{(Career orientation)}}
\begin{itemize}
\item \cdkv{Lập luận}{Nhiều học sinh chọn sai ngành đại học. Một năm nghỉ ngơi giúp họ suy nghĩ kỹ mình thích gì.}
\item \cdkv{Dẫn chứng/Giải thích}{Thay vì lãng phí 4 năm học sai ngành, 1 năm trải nghiệm thực tế giúp họ chọn đúng nghề nghiệp phù hợp với sở thích.}
\end{itemize}

Dàn ý tham khảo (Sample Outline)
\begin{itemize}
\item \cdkv{Mở bài}{Giới thiệu chủ đề Gap Year. Đưa ra quan điểm: Lợi ích lớn hơn tác hại.}
\item \cdblue{\textbf{Thân bài 1 (Tác hại):}}
\begin{itemize}
\item \cdkv{Câu chủ đề}{Có vài bất lợi khi nghỉ 1 năm.}
\item \cdkv{Ý 1}{Có thể mất động lực học.}
\item \cdkv{Ý 2}{Tốt nghiệp muộn hơn bạn bè.}
\end{itemize}
\item \cdblue{\textbf{Thân bài 2 (Lợi ích - Quan trọng hơn):}}
\begin{itemize}
\item \cdkv{Câu chủ đề}{Tuy nhiên, lợi ích mang lại là rất lớn.}
\item \cdkv{Ý 1}{Có thêm kinh nghiệm thực tế và kỹ năng mềm.}
\item \cdkv{Ý 2}{Hiểu rõ bản thân muốn gì, tránh chọn sai ngành.}
\end{itemize}
\item \cdkv{Kết bài}{Khẳng định lại: Dù có rủi ro, nhưng việc trải nghiệm thực tế là rất đáng giá.}
\end{itemize}

\cdsection{Phần II: Kiến tạo bài luận hoàn chỉnh}

Dưới đây là bài mẫu viết theo phong cách đơn giản, dễ học.

\cdstep{1. Mở bài \textit{(Introduction)}}
\begin{itemize}
\item \cdkv{Nhiệm vụ}{Paraphrase (viết lại) đề bài và trả lời câu hỏi.}
\item \cdblue{\textbf{Ví dụ Mẫu:}}
\end{itemize}

Nowadays, some students choose to work or travel for a year before going to university. In my opinion, although there are some disadvantages, I believe the advantages of this trend are greater.

(Ngày nay, một số học sinh chọn làm việc hoặc du lịch 1 năm trước khi vào đại học. Theo ý kiến của tôi, mặc dù có vài bất lợi, tôi tin rằng lợi ích của xu hướng này lớn hơn.)

\cdstep{2. Thân bài 1 \textit{(Body Paragraph 1)}: Bàn về Tác hại}
\begin{itemize}
\item \cdkv{Nhiệm vụ}{Nêu ra 1-2 điểm xấu (viết đơn giản).}
\item \cdkv{Câu chủ đề}{On the one hand, taking a gap year can bring some problems.}
\end{itemize}

(Một mặt, nghỉ một năm có thể mang lại vài vấn đề.)
\begin{itemize}
\item \cdblue{\textbf{Giải thích \& Ví dụ:}}
\end{itemize}

First, students might lose their study habits. After working for a long time, returning to school can be difficult because they are used to earning money. Second, they will graduate later than their friends, which might make them feel stressed.

(Đầu tiên, học sinh có thể mất thói quen học tập. Sau khi đi làm một thời gian dài, quay lại trường có thể khó khăn vì họ đã quen kiếm tiền. Thứ hai, họ sẽ tốt nghiệp muộn hơn bạn bè, điều này có thể làm họ thấy áp lực.)

\cdstep{3. Thân bài 2 \textit{(Body Paragraph 2)}: Bàn về Lợi ích (Quan trọng hơn)}
\begin{itemize}
\item \cdkv{Nhiệm vụ}{Nhấn mạnh vào những điều tốt đẹp (Kinh nghiệm, kỹ năng).}
\item \cdkv{Câu chủ đề}{On the other hand, I think the benefits of working or traveling are more significant.}
\end{itemize}

(Mặt khác, tôi nghĩ lợi ích của việc đi làm hoặc du lịch là đáng kể hơn.)
\begin{itemize}
\item \cdblue{\textbf{Giải thích \& Ví dụ:}}
\end{itemize}

Firstly, young people can learn many soft skills. For example, if they have a job, they learn how to communicate and manage money. These skills are very useful for their future life. Secondly, a gap year gives them time to think about their career. This helps them choose the right major in university.

(Đầu tiên, người trẻ có thể học nhiều kỹ năng mềm. Ví dụ, nếu họ có một công việc, họ học cách giao tiếp và quản lý tiền bạc. Những kỹ năng này rất hữu ích cho cuộc sống sau này. Thứ hai, một năm nghỉ ngơi cho họ thời gian suy nghĩ về sự nghiệp. Điều này giúp họ chọn đúng chuyên ngành ở đại học.)

\cdstep{4. Kết bài \textit{(Conclusion)}}
\begin{itemize}
\item \cdkv{Nhiệm vụ}{Tóm tắt lại ý chính.}
\item \cdblue{\textbf{Ví dụ Mẫu:}}
\end{itemize}

In conclusion, taking a gap year has some drawbacks regarding study habits. However, I think it is a good idea because students can gain valuable experience and skills. The advantages clearly outweigh the disadvantages.

(Tóm lại, nghỉ một năm có vài hạn chế về thói quen học tập. Tuy nhiên, tôi nghĩ đó là ý hay vì học sinh có thể đạt được kinh nghiệm và kỹ năng quý giá. Lợi ích rõ ràng vượt trội hơn tác hại.)

\cdsection{Phần III: Công cụ ngôn ngữ}

Dưới đây là các từ vựng từ đơn giản đến trung bình, dễ nhớ và dễ áp dụng cho bài viết này.

Từ/ Cụm từ

Nghĩa tiếng Việt

Ví dụ đơn giản (Example)

Gap year (n)

Năm nghỉ phép (trước đại học)

Taking a gap year is becoming popular.

Drawback (n)

Mặt hạn chế, bất lợi (= Disadvantage)

One drawback is that you graduate late.

Benefit (n)

Lợi ích (= Advantage)

The main benefit is earning money.

Soft skills (n)

Kỹ năng mềm (giao tiếp, làm việc nhóm)

Schools do not always teach soft skills.

Mature (adj)

Trưởng thành

Working helps students become more mature.

Independent (adj)

Tự lập

Living away from home makes you independent.

Valuable experience

Kinh nghiệm quý giá

Traveling gives us valuable experience.

Lose momentum

Mất đà, mất hứng thú

Students might lose momentum in studying.

Broaden horizons

Mở rộng tầm mắt (hiểu biết)

Traveling helps to broaden horizons.

Tuition fees (n)

Học phí

They work to save money for tuition fees.

\end{document}
